\subsection{\texorpdfstring{第~\ref{sec:affine-connections}~节}{仿射联络}}\label{subsec:solutions-affine-connections}

\begin{description}[font=\normalfont,leftmargin=2em,labelwidth=1.5em]
\RGsolution{ex:locality-and-coordinate-expressions}{locality-and-coordinate-expressions}
\begin{enumerate}[label=(\alph*),ref=\alph*,leftmargin=*]
\item 设 $p\in W$, $\rho\in C^\infty(M)$ 是满足 $\operatorname{supp}\rho\subset W$ 和 $\rho(p)>0$ 的隆起函数. 则
\[
  \rho X=\rho\widetilde X\quad\text{且}\quad \rho Y=\rho\widetilde Y.
\]
由仿射联络的性质, 得
\[
  \nabla_{\rho X}(\rho Y)=\rho(\rho\nabla_XY+(X\cdot\rho)Y),
\]
类似地,
\[
  \nabla_{\rho\widetilde X}(\rho\widetilde Y)
  =\rho(\rho\nabla_{\widetilde X}\widetilde Y+(\widetilde X\cdot\rho)\widetilde Y).
\]
令两个表达式相等, 得
\[
  \rho\nabla_XY+(X\cdot\rho)Y
  =\rho\nabla_{\widetilde X}\widetilde Y+(\widetilde X\cdot\rho)\widetilde Y,
\]
在 $p$ 处即为
\[
  \rho(p)(\nabla_XY)_p+(X_p\cdot\rho)Y_p
  =\rho(p)(\nabla_{\widetilde X}\widetilde Y)_p+(\widetilde X_p\cdot\rho)\widetilde Y_p.
\]
由于 $p\in W$, 有 $X_p=\widetilde X_p$ 且 $Y_p=\widetilde Y_p$. 因此,
\[
  \rho(p)(\nabla_XY)_p=\rho(p)(\nabla_{\widetilde X}\widetilde Y)_p
  \iff(\nabla_XY)_p=(\nabla_{\widetilde X}\widetilde Y)_p,
\]
其中使用了 $\rho(p)>0$. 因为 $p\in W$ 是任取的, 所以 $\nabla_XY=\nabla_{\widetilde X}\widetilde Y$ 在 $W$ 上成立.
\item 有
\begin{align*}
  \nabla_XY&=\nabla_X\left(\sum_{i=1}^nY^i\frac{\partial}{\partial x^i}\right)
  =\sum_{i=1}^n(X\cdot Y^i)\frac{\partial}{\partial x^i}
  +\sum_{i=1}^nY^i\nabla_X\frac{\partial}{\partial x^i}\\
  &=\sum_{i=1}^n(X\cdot Y^i)\frac{\partial}{\partial x^i}
  +\sum_{k=1}^nY^k\nabla_{\left(\sum_{j=1}^n X^j\frac{\partial}{\partial x^j}\right)}
       \frac{\partial}{\partial x^k}\\
  &=\sum_{i=1}^n(X\cdot Y^i)\frac{\partial}{\partial x^i}
  +\sum_{j,k=1}^nX^jY^k\nabla_{\frac{\partial}{\partial x^j}}\frac{\partial}{\partial x^k}\\
  &=\sum_{i=1}^n(X\cdot Y^i)\frac{\partial}{\partial x^i}
  +\sum_{i,j,k=1}^nX^jY^k\Gamma^i_{jk}\frac{\partial}{\partial x^i}.
\end{align*}
\item 利用
\[
  \dot c(t)=\sum_{i=1}^n\dot{x}^i(t)\left(\frac{\partial}{\partial x^i}\right)_{c(t)}
  \quad\text{和}\quad
  V(t)=\sum_{i=1}^n V^i(t)\left(\frac{\partial}{\partial x^i}\right)_{c(t)},
\]
得到
\begin{align*}
  \frac{DV}{\mathrm{d}t}(t)&=\nabla_{\dot c(t)}V\\
  &=\sum_{i=1}^n\left(\dot c(t)\cdot V^i(t)+\sum_{j,k=1}^n\Gamma^i_{jk}(c(t))\dot{x}^j(t)V^k(t)\right)
           \left(\frac{\partial}{\partial x^i}\right)_{c(t)}\\
  &=\sum_{i=1}^n\left(\dot V^i(t)+\sum_{j,k=1}^n\Gamma^i_{jk}(c(t))\dot{x}^j(t)V^k(t)\right)
           \left(\frac{\partial}{\partial x^i}\right)_{c(t)}.
\end{align*}
因此, 平行移动定律的坐标方程为
\[
  \dot V^i(t)+\sum_{j,k=1}^n\Gamma^i_{jk}(c(t))\dot{x}^j(t)V^k(t)=0
  \quad(i=1,\ldots,n).
\]
\item 在 $V=\dot c$, 即 $V^i(t)=\dot{x}^i(t)$ 的情形下应用 \eqref{eq:parallel-transport-in-local-coordinates}, 得
\[
  \ddot{x}^i(t)+\sum_{j,k=1}^n\Gamma^i_{jk}(c(t))\dot{x}^j(t)\dot{x}^k(t)=0
  \quad(i=1,\ldots,n).
\]
\end{enumerate}

\RGsolution{ex:covariant-derivative-of-tensors}{covariant-derivative-of-tensors}
\begin{enumerate}[label=(\alph*),ref=\alph*,leftmargin=*]
\item 设 $X,Y,Z\in\mathfrak X(M)$ 且 $f,g\in C^\infty(M)$. 则
\begin{align*}
  (\nabla_X\omega)(fY+gZ)
  &=X\cdot(\omega(fY+gZ))-\omega(\nabla_X(fY+gZ))\\
  &=X\cdot(f\omega(Y)+g\omega(Z))\\
  &\quad-\omega(f\nabla_XY+(X\cdot f)Y+g\nabla_XZ+(X\cdot g)Z)\\
  &=(X\cdot f)\omega(Y)+(X\cdot g)\omega(Z)+f(X\cdot(\omega(Y))-\omega(\nabla_XY))\\
  &\quad+g(X\cdot(\omega(Z))-\omega(\nabla_XZ))-(X\cdot f)\omega(Y)-(X\cdot g)\omega(Z)\\
  &=(f\nabla_X\omega)(Y)+g(\nabla_X\omega)(Z).
\end{align*}
\item
\begin{enumerate}[label=(\roman*),ref=\roman*,leftmargin=*]
\item 有
\begin{align*}
  (\nabla_{fX+gY}\omega)(Z)&=(fX+gY)\cdot(\omega(Z))-\omega(\nabla_{fX+gY}Z)\\
  &=fX\cdot(\omega(Z))+gY\cdot(\omega(Z))-f\omega(\nabla_X Z)-g\omega(\nabla_Y Z)\\
  &=f(\nabla_X\omega+g\nabla_Y\omega)(Z);
\end{align*}
\item 此外,
\begin{align*}
  (\nabla_X(\omega+\eta))(Y)
  &=X\cdot((\omega+\eta)(Y))-(\omega+\eta)(\nabla_XY)\\
  &=X\cdot(\omega(Y))+X\cdot(\eta(Y))-\omega(\nabla_XY)-\eta(\nabla_XY)\\
  &=(\nabla_X\omega+\nabla_X\eta)(Y);
\end{align*}
\item 最后,
\begin{align*}
  (\nabla_X(f\omega))(Y)
  &=X\cdot(f\omega(Y))-f\omega(\nabla_XY)\\
  &=(X\cdot f)\omega(Y)+f(X\cdot(\omega(Y))-\omega(\nabla_XY))\\
  &=((X\cdot f)\omega+f\nabla_X\omega)(Y).
\end{align*}
\end{enumerate}
\item 在这些坐标下, 有
\begin{align*}
  \nabla_X\omega(Y)
  &=X\cdot(\omega(Y))-\omega(\nabla_XY)\\
  &=X\cdot\left(\sum_{i=1}^n\omega_i\,\mathrm{d}x^i(Y)\right)
     -\omega\left(\sum_{i=1}^n\left(X\cdot Y^i+\sum_{j,k=1}^n\Gamma^i_{jk}X^jY^k\right)
         \frac{\partial}{\partial x^i}\right)\\
  &=X\cdot\left(\sum_{i=1}^n\omega_iY^i\right)-\sum_{i=1}^n\omega_i(X\cdot Y^i)
      -\sum_{i,j,k=1}^n\Gamma^i_{jk}X^jY^k\omega_i\\
  &=\sum_{i=1}^n(X\cdot\omega_i)Y^i-\sum_{i=1}^n\left(\sum_{j,k=1}^n\Gamma^k_{ji}X^j\omega_k\right)Y^i\\
  &=\left(\sum_{i=1}^n\left(X\cdot\omega_i-\sum_{j,k=1}^n\Gamma^k_{ji}X^j\omega_k\right)\mathrm{d}x^i\right)(Y).
\end{align*}
\item 设 $X$ 是向量场, $T$ 是 $(k,m)$-张量场. 对所有 $X_1,\ldots,X_k\in\mathfrak X(M)$ 和 $\theta^1,\ldots,\theta^m\in\Omega^1(M)$, 定义
\begin{align*}
  &\nabla_XT(X_1,\ldots,X_k,\theta^1,\ldots,\theta^m)
    :=X\cdot(T(X_1,\ldots,X_k,\theta^1,\ldots,\theta^m))\\
  &\quad-T(\nabla_XX_1,\ldots,X_k,\theta^1,\ldots,\theta^m)
      -\cdots-T(X_1,\ldots,\nabla_XX_k,\theta^1,\ldots,\theta^m)\\
  &\quad-T(X_1,\ldots,X_k,\nabla_X\theta^1,\ldots,\theta^m)
      -\cdots-T(X_1,\ldots,X_k,\theta^1,\ldots,\nabla_X\theta^m).
\end{align*}
注意, 这推广了 1-形式的定义, 且当 $T$ 为向量场时与通常的定义一致. 与 1-形式情形类似的计算给出
\[
  \begin{aligned}
  \nabla_XT={}&\sum_{i_1,\ldots,i_k,j_1,\ldots,j_m=1}^n
  \Biggl( X\cdot T_{i_1\cdots i_k}^{j_1\cdots j_m}
         -\sum_{r,s=1}^n\Gamma^s_{ri_1}X^rT_{si_2\cdots i_k}^{j_1\cdots j_m}-\cdots\\
  &\quad-\sum_{r,s=1}^n\Gamma^s_{ri_k}X^rT_{i_1\cdots i_{k-1}s}^{j_1\cdots j_m}
        +\sum_{r,s=1}^n\Gamma^{j_1}_{rs}X^rT_{i_1\cdots i_k}^{sj_2\cdots j_m}+\cdots\\
  &\quad+\sum_{r,s=1}^n\Gamma^{j_m}_{rs}X^rT_{i_1\cdots i_k}^{j_1\cdots j_{k-1}s}\Biggr)
    \mathrm{d}x^{i_1}\otimes\cdots\otimes\mathrm{d}x^{i_k}
    \otimes\frac{\partial}{\partial x^{j_1}}\otimes\cdots\otimes\frac{\partial}{\partial x^{j_m}}.
  \end{aligned}
\]
\end{enumerate}
\end{description}
